% Options for packages loaded elsewhere
\PassOptionsToPackage{unicode}{hyperref}
\PassOptionsToPackage{hyphens}{url}
%
\documentclass[
]{article}
\usepackage{amsmath,amssymb}
\usepackage{lmodern}
\usepackage{iftex}
\ifPDFTeX
  \usepackage[T1]{fontenc}
  \usepackage[utf8]{inputenc}
  \usepackage{textcomp} % provide euro and other symbols
\else % if luatex or xetex
  \usepackage{unicode-math}
  \defaultfontfeatures{Scale=MatchLowercase}
  \defaultfontfeatures[\rmfamily]{Ligatures=TeX,Scale=1}
\fi
% Use upquote if available, for straight quotes in verbatim environments
\IfFileExists{upquote.sty}{\usepackage{upquote}}{}
\IfFileExists{microtype.sty}{% use microtype if available
  \usepackage[]{microtype}
  \UseMicrotypeSet[protrusion]{basicmath} % disable protrusion for tt fonts
}{}
\makeatletter
\@ifundefined{KOMAClassName}{% if non-KOMA class
  \IfFileExists{parskip.sty}{%
    \usepackage{parskip}
  }{% else
    \setlength{\parindent}{0pt}
    \setlength{\parskip}{6pt plus 2pt minus 1pt}}
}{% if KOMA class
  \KOMAoptions{parskip=half}}
\makeatother
\usepackage{xcolor}
\IfFileExists{xurl.sty}{\usepackage{xurl}}{} % add URL line breaks if available
\IfFileExists{bookmark.sty}{\usepackage{bookmark}}{\usepackage{hyperref}}
\hypersetup{
  pdftitle={Homework 2: DAGs and E-Values - Solution},
  pdfauthor={Xuange Liang},
  hidelinks,
  pdfcreator={LaTeX via pandoc}}
\urlstyle{same} % disable monospaced font for URLs
\usepackage{longtable,booktabs,array}
\usepackage{calc} % for calculating minipage widths
% Correct order of tables after \paragraph or \subparagraph
\usepackage{etoolbox}
\makeatletter
\patchcmd\longtable{\par}{\if@noskipsec\mbox{}\fi\par}{}{}
\makeatother
% Allow footnotes in longtable head/foot
\IfFileExists{footnotehyper.sty}{\usepackage{footnotehyper}}{\usepackage{footnote}}
\makesavenoteenv{longtable}
\usepackage{graphicx}
\makeatletter
\def\maxwidth{\ifdim\Gin@nat@width>\linewidth\linewidth\else\Gin@nat@width\fi}
\def\maxheight{\ifdim\Gin@nat@height>\textheight\textheight\else\Gin@nat@height\fi}
\makeatother
% Scale images if necessary, so that they will not overflow the page
% margins by default, and it is still possible to overwrite the defaults
% using explicit options in \includegraphics[width, height, ...]{}
\setkeys{Gin}{width=\maxwidth,height=\maxheight,keepaspectratio}
% Set default figure placement to htbp
\makeatletter
\def\fps@figure{htbp}
\makeatother
\setlength{\emergencystretch}{3em} % prevent overfull lines
\providecommand{\tightlist}{%
  \setlength{\itemsep}{0pt}\setlength{\parskip}{0pt}}
\setcounter{secnumdepth}{5}
% Shared LaTeX style for homework submissions
\usepackage{microtype}
\usepackage{booktabs}
\usepackage{longtable}
\usepackage{array}
\usepackage{caption}
\usepackage{fancyhdr}
\usepackage{titlesec}
\usepackage{setspace}
\usepackage{float}
\usepackage{graphicx}
\usepackage[table]{xcolor}
\usepackage{enumitem}
\usepackage[most]{tcolorbox}

\definecolor{TitleBlue}{HTML}{0B3C5D}
\definecolor{SoftGray}{HTML}{F6F8FA}

\setstretch{1.15}
\setlength{\parskip}{4pt}
\setlength{\parindent}{0pt}

\titleformat{\section}{\large\bfseries\color{TitleBlue}}{\thesection}{0.6em}{}
\titleformat{\subsection}{\normalsize\bfseries\color{TitleBlue}}{\thesubsection}{0.6em}{}

\captionsetup{
  font=small,
  labelfont=bf
}

\rowcolors{2}{SoftGray}{white}
\setlist[itemize]{leftmargin=*, topsep=2pt, itemsep=1pt}
\setlist[enumerate]{leftmargin=*, topsep=2pt, itemsep=1pt}

\pagestyle{fancy}
\fancyhf{}
\fancyhead[L]{\nouppercase{\leftmark}}
\fancyhead[R]{P8586 Applied Methods}
\fancyfoot[C]{\thepage}
\setlength{\headheight}{14pt}

\AtBeginDocument{
  \hypersetup{
    colorlinks=true,
    linkcolor=TitleBlue,
    urlcolor=TitleBlue,
    citecolor=TitleBlue
  }
}

% Keep title and TOC on same page, content starts on new page
\AtBeginDocument{
  \let\oldtableofcontents\tableofcontents
  \renewcommand{\tableofcontents}{%
    \oldtableofcontents
    \clearpage
  }
}
\ifLuaTeX
  \usepackage{selnolig}  % disable illegal ligatures
\fi

\title{Homework 2: DAGs and E-Values - Solution}
\author{Xuange Liang}
\date{February 12, 2026}

\begin{document}
\maketitle

{
\setcounter{tocdepth}{3}
\tableofcontents
}
\hypertarget{literature-review}{%
\section{Literature Review}\label{literature-review}}

\textbf{Perform a literature review on variables/factors confounding the
association between IVH and mortality.}

Based on the literature review, several key risk factors were identified
that are associated with both Intraventricular Hemorrhage (IVH) and
Neonatal Mortality in Very Low Birth Weight (VLBW) infants.

\begin{longtable}[]{@{}
  >{\raggedright\arraybackslash}p{(\columnwidth - 10\tabcolsep) * \real{0.17}}
  >{\raggedright\arraybackslash}p{(\columnwidth - 10\tabcolsep) * \real{0.17}}
  >{\raggedright\arraybackslash}p{(\columnwidth - 10\tabcolsep) * \real{0.17}}
  >{\raggedright\arraybackslash}p{(\columnwidth - 10\tabcolsep) * \real{0.17}}
  >{\raggedright\arraybackslash}p{(\columnwidth - 10\tabcolsep) * \real{0.17}}
  >{\raggedright\arraybackslash}p{(\columnwidth - 10\tabcolsep) * \real{0.17}}@{}}
\toprule
Publications (Example Sources) & Potential Confounding Factors &
Association w/ IVH & Estimates (95\% CI) & Association w/ Mortality &
Estimates (95\% CI) \\
\midrule
\endhead
\href{https://doi.org/10.1038/jp.2016.38}{\emph{Wei et al.~(2016)}} /
\href{https://doi.org/10.1002/14651858.CD004454.pub3}{\emph{Roberts et
al.~(2017)}} & \textbf{Antenatal Corticosteroids (ACS)} & Protective
(Reduces IVH risk) & Severe IVH: OR = 0.51 (0.45-0.58)
\href{https://doi.org/10.1038/jp.2016.38}{{[}1{]}} & Protective (Reduces
Mortality) & Neonatal death: RR = 0.69 (0.58-0.81)
\href{https://doi.org/10.1002/14651858.CD004454.pub3}{{[}2{]}} \\
\href{https://doi.org/10.3389/fphys.2018.01253}{\emph{Villamor-Martinez
et al.~(2018)}} & \textbf{Chorioamnionitis (Clinical/Histological)} &
Risk Factor (Increases IVH risk) & OR = 1.88 (1.61-2.19)
\href{https://doi.org/10.3389/fphys.2018.01253}{{[}3{]}} & Associated
with prematurity/sepsis pathways & (Estimates vary by outcome
definition/study) \\
\emph{HW1 Data / Literature} & \textbf{Gestational Age} & Strong
Negative Assoc. & (See HW1 results) & Strong Negative Assoc. & (See HW1
results) \\
\emph{HW1 Data / Literature} & \textbf{Birth Weight} & Strong Negative
Assoc. & (See HW1 results) & Strong Negative Assoc. & (See HW1
results) \\
\href{https://doi.org/10.1186/s12884-022-05245-2}{\emph{Zhao et
al.~(2022)}} /
\href{https://doi.org/10.1111/j.1651-2227.2010.01935.x}{\emph{Lee et
al.~(2010)}} & \textbf{Low 5-min Apgar score} & Risk Factor (Increases
IVH risk) & Apgar \textless=7: OR = 2.273 (1.163-4.442)
\href{https://doi.org/10.1186/s12884-022-05245-2}{{[}4{]}} & Risk Factor
(Increases Mortality) & Example (24 wks): RR = 3.1 (2.9-3.4) for Apgar
0-3 vs 7-10
\href{https://doi.org/10.1111/j.1651-2227.2010.01935.x}{{[}5{]}} \\
\bottomrule
\end{longtable}

\begin{center}\rule{0.5\linewidth}{0.5pt}\end{center}

\hypertarget{directed-acyclic-graph-dag}{%
\section{Directed Acyclic Graph
(DAG)}\label{directed-acyclic-graph-dag}}

\textbf{Draw a Directed Acyclic Graph (DAG) to illustrate the conceptual
framework of the causal association between IVH and infant death.}

Below is the DAG written in \textbf{DAGitty syntax} (copy/paste into
\url{https://dagitty.net/dags.html}):

\begin{figure}
\centering
\includegraphics[width=0.92\textwidth,height=\textheight]{DAG_HW2.png}
\caption{Final DAG for IVH and mortality}
\end{figure}

\newpage

\textbf{DAGitty Code:}

\begin{tcolorbox}[
  colback=SoftGray,
  colframe=TitleBlue,
  boxrule=0.5pt,
  arc=2pt,
  left=6pt,right=6pt,top=6pt,bottom=6pt,
  before skip=4pt, after skip=10pt
]
\begingroup\footnotesize
\noindent\begin{minipage}[t]{0.49\linewidth}
\begin{verbatim}
dag {
bb="-1,-1,9,6"

IVH [exposure,pos="4,3"]
DEAD [outcome,pos="8,3"]
GEST_AGE [pos="2,4"]
BWT [pos="3,5"]
ACS [pos="1,5"]
CHORIO [pos="0,4"]
DELIVERY [pos="2,2"]
APGAR5 [pos="3,2"]
SEX [pos="2,1"]
PNEUMO [pos="6,2"]

IVH -> DEAD
IVH -> PNEUMO
PNEUMO -> DEAD

GEST_AGE -> IVH
GEST_AGE -> DEAD
GEST_AGE -> BWT
GEST_AGE -> APGAR5
GEST_AGE -> PNEUMO

BWT -> IVH
BWT -> DEAD
\end{verbatim}
\end{minipage}\hfill
\begin{minipage}[t]{0.49\linewidth}
\begin{verbatim}
ACS -> IVH
ACS -> DEAD

CHORIO -> GEST_AGE
CHORIO -> IVH
CHORIO -> DEAD

DELIVERY -> IVH
DELIVERY -> DEAD

APGAR5 -> IVH
APGAR5 -> DEAD

SEX -> IVH
SEX -> DEAD
}
\end{verbatim}
\end{minipage}
\endgroup
\end{tcolorbox}

\textbf{DAGitty analysis (for total effect of IVH on death):} - Minimal
sufficient adjustment set identified by this DAG:\\
\texttt{\{GEST\_AGE,\ BWT,\ ACS,\ CHORIO,\ DELIVERY,\ APGAR5,\ SEX\}} -
\texttt{PNEUMO} is modeled as a post-exposure mediator
(\texttt{IVH\ -\textgreater{}\ PNEUMO\ -\textgreater{}\ DEAD}), so it
should \textbf{not} be adjusted for when estimating the total causal
effect of IVH.

\begin{center}\rule{0.5\linewidth}{0.5pt}\end{center}

\hypertarget{comparison-with-hw1-model}{%
\section{Comparison with HW1 Model}\label{comparison-with-hw1-model}}

\textbf{Compare variables identified in the DAG to the covariates
included in the adjusted regression model in Homework 1.}

\begin{enumerate}
\def\labelenumi{\arabic{enumi}.}
\item
  \textbf{Covariates in HW1 adjusted model:} \texttt{pneumo},
  \texttt{bwt\_cat}, \texttt{gest\_cat}.
\item
  \textbf{DAGitty-based target adjustment set (if sample size is not
  limiting):}\\
  \texttt{GEST\_AGE}, \texttt{BWT}, \texttt{ACS}, \texttt{CHORIO},
  \texttt{DELIVERY}, \texttt{APGAR5}, \texttt{SEX}.
\item
  \textbf{Variables missing from HW1 (relative to DAGitty set):}

  \begin{itemize}
  \tightlist
  \item
    \texttt{ACS} (antenatal corticosteroids)
  \item
    \texttt{CHORIO} (chorioamnionitis)
  \item
    \texttt{DELIVERY} (delivery mode)
  \item
    \texttt{APGAR5} (5-min Apgar)
  \item
    \texttt{SEX} (infant sex)
  \end{itemize}
\item
  \textbf{Role of pneumothorax in causal modeling:}

  \begin{itemize}
  \tightlist
  \item
    In this DAG, pneumothorax is primarily post-exposure and on the
    pathway
    \texttt{IVH\ -\textgreater{}\ PNEUMO\ -\textgreater{}\ DEAD}.
  \item
    Therefore, for estimating the \textbf{total effect} of IVH on
    mortality, pneumothorax should not be adjusted for.
  \item
    If the research target were a controlled direct effect (not total
    effect), mediation-oriented methods would be required.
  \end{itemize}
\item
  \textbf{Conclusion for HW2 Question 3:}

  \begin{itemize}
  \tightlist
  \item
    Yes, additional confounders should have been included if data and
    sample size allowed.
  \item
    I would add back \texttt{ACS}, \texttt{CHORIO}, \texttt{DELIVERY},
    \texttt{APGAR5}, and \texttt{SEX}, and avoid routine adjustment for
    \texttt{pneumo} in the primary total-effect model.
  \end{itemize}
\end{enumerate}

\begin{center}\rule{0.5\linewidth}{0.5pt}\end{center}

\hypertarget{e-values-calculation}{%
\section{E-Values Calculation}\label{e-values-calculation}}

\textbf{Calculate the E-values for the exposure of interest (IVH). Do
you think the association is likely to be explained away by unmeasured
confounders?}

\textbf{Data from Homework 1 (Poisson Model):} * \textbf{Adjusted Risk
Ratio (\(RR_{obs}\))}: 2.059 (from HW1 final Poisson model) *
\textbf{95\% Confidence Interval}: 1.351 - 3.137

\textbf{Formula:} \[ E\text{-value} = RR + \sqrt{RR \times (RR - 1)} \]

\textbf{Computed with R (\texttt{EValue} package; confirmed by manual
formula):}

\begin{enumerate}
\def\labelenumi{\arabic{enumi}.}
\item
  \textbf{E-value for point estimate} (\(RR=2.059\)):\\
  \[ \mathbf{E = 3.536} \]
\item
  \textbf{E-value for lower confidence limit} (\(LL=1.351\)):\\
  \[ \mathbf{E_{LL} = 2.040} \]
\end{enumerate}

\textbf{Interpretation:} * To explain away the observed risk ratio of
2.059, an unmeasured confounder would need to have an association with
both IVH and death of at least \textbf{RR = 3.536}, conditional on
measured covariates. * To move the lower CI bound to the null, an
unmeasured confounder would still need an association of at least
\textbf{RR = 2.040} with both exposure and outcome.

\textbf{Conclusion:} * An E-value of 3.536 is relatively large in this
clinical context. Few single unmeasured confounders (e.g.,
chorioamnionitis with effect sizes often around 1.9-2.2) are likely to
be strong enough on both links after adjustment for prematurity-related
factors. * Therefore, it is \textbf{unlikely} that the observed
association between IVH and infant mortality is entirely explained away
by unmeasured confounding. The evidence supports IVH as an independent
risk factor.

\begin{center}\rule{0.5\linewidth}{0.5pt}\end{center}

\hypertarget{references}{%
\section{References}\label{references}}

\begin{enumerate}
\def\labelenumi{\arabic{enumi}.}
\tightlist
\item
  Wei JC, Catalano R, Profit J, Gould JB, Lee HC. Impact of antenatal
  steroids on intraventricular hemorrhage in very-low-birth weight
  infants. \emph{J Perinatol}. 2016;36(5):352-356. doi:
  \href{https://doi.org/10.1038/jp.2016.38}{10.1038/jp.2016.38}
\item
  Roberts D, Brown J, Medley N, Dalziel SR. Antenatal corticosteroids
  for accelerating fetal lung maturation for women at risk of preterm
  birth. \emph{Cochrane Database Syst Rev}. 2017;3(3):CD004454. doi:
  \href{https://doi.org/10.1002/14651858.CD004454.pub3}{10.1002/14651858.CD004454.pub3}
\item
  Villamor-Martinez E, et al.~Chorioamnionitis is a risk factor for
  intraventricular hemorrhage in preterm infants: a systematic review
  and meta-analysis. \emph{Front Physiol}. 2018;9:1253. doi:
  \href{https://doi.org/10.3389/fphys.2018.01253}{10.3389/fphys.2018.01253}
\item
  Zhao Y, Zhang W, Tian X. Analysis of risk factors of early
  intraventricular hemorrhage in very-low-birth-weight premature
  infants: a single center retrospective study. \emph{BMC Pregnancy
  Childbirth}. 2022;22:890. doi:
  \href{https://doi.org/10.1186/s12884-022-05245-2}{10.1186/s12884-022-05245-2}
\item
  Lee HC, Subeh M, Gould JB. Low Apgar score and mortality in extremely
  preterm neonates born in the United States. \emph{Acta Paediatr}.
  2010;99(12):1785-1789. doi:
  \href{https://doi.org/10.1111/j.1651-2227.2010.01935.x}{10.1111/j.1651-2227.2010.01935.x}
\end{enumerate}

\end{document}
